\documentclass[11pt]{article}
\usepackage[a4paper,margin=1.1in]{geometry}
\usepackage[T1]{fontenc}
\usepackage{lmodern}
\usepackage{amsmath,amssymb,amsthm}
\usepackage{booktabs}
\usepackage{longtable}
\usepackage{array}
\usepackage{microtype}
\usepackage[hidelinks]{hyperref}
\usepackage{natbib}
\usepackage{setspace}
\usepackage{enumitem}
\usepackage{xcolor}

\newtheorem{theorem}{Theorem}[section]
\newtheorem{proposition}[theorem]{Proposition}
\newtheorem{lemma}[theorem]{Lemma}
\newtheorem{corollary}[theorem]{Corollary}
\theoremstyle{definition}
\newtheorem{definition}[theorem]{Definition}
\newtheorem{assumption}[theorem]{Assumption}
\theoremstyle{remark}
\newtheorem{remark}[theorem]{Remark}

% A formal name, as it appears in the Lean development.
\newcommand{\leanunderscore}{\char`\_\hskip 0pt plus 0.4pt\relax}
\newcommand{\lean}[1]{\begingroup\ttfamily\small\let\_\leanunderscore #1\endgroup}
% The Lean name attached to a stated result.
\newcommand{\formal}[1]{\par\noindent\begingroup\footnotesize\raggedright\emph{Formal statement:} \lean{#1}.\par\endgroup}
\newcommand{\Thorn}{\text{\th}}
\newcommand{\R}{\mathbb{R}}
\newcommand{\E}{\mathbb{E}}
\DeclareMathOperator{\supp}{supp}

\onehalfspacing
\tolerance=600 \emergencystretch=1.5em

\title{Overlapping Generations in Lean:\\ The Life-Cycle Household}
\author{LeanEconomics Org [mostly Claude]\thanks{The Lean development described here is at
\url{https://github.com/LeanEconomics/LeanEconomics}. It was written mostly by Claude (Anthropic),
an AI coding model, working under the direction of Robert Kirkby (Victoria University of
Wellington); every result reported was checked by the Lean kernel. This is the companion to the
write-up on Aiyagari (1994) in the same repository. Comments welcome.}}
\date{September 2026\\[0.5em]\small Draft, in progress}

\begin{document}
\maketitle

\begin{abstract}
\noindent
This is the record of a programme to formalise overlapping-generations economies in Lean 4, and
of its first stage: the life-cycle household. A household that lives $J$ periods against
uninsurable earnings risk and a borrowing constraint solves a finite backward induction, and
everything proved for the infinite-horizon Bellman operator with an arbitrary continuation applies
stage by stage with no fixed point. Beyond that we prove the finite-horizon sandwich: with $k$
periods still to come, consumption lies between $\kappa_k m$ and $\kappa_k(m+H_k(z))$, where
$1/\kappa_k=\sum_{i\le k}(\Thorn/R)^i$ is the consumption share of a household with no future
income and $H_k$ is the expected present value of the next $k$ incomes; the propensities decrease
in the horizon towards the infinite-horizon $\kappa=1-\Thorn/R$. We then add retirement and a
pension. With a lump-sum pension the retired phase is an income-fluctuation problem in its own
right and everything survives. With a means-tested pension we settle a claim often made
informally: a means test does make the policy non-monotone, but not for the reason usually given.
Optimal saving rises with cash on hand whatever the continuation, by Topkis, so a value function
made non-concave by the test cannot by itself reverse it; what the test reverses is consumption,
which jumps down at the wealth where the household gives the pension up, and cash on hand itself,
once the taper is steeper than the gross return. The programme's target is
uniqueness of the stationary equilibrium at relative risk aversion above one, which the infinite
horizon could not close: in a life cycle every comparative static is a finite induction, wealth is
bounded, and the age-one distribution is common to every rate.
\end{abstract}

\noindent\emph{JEL:} C63, D52, E21. \quad \emph{Keywords:} overlapping generations, life cycle,
Bewley models, formal verification, Lean.

\section{Introduction}

The companion write-up in this repository proved that the stationary equilibrium of
\citet{Aiyagari1994} is unique at his logarithmic calibration, and located exactly what stands in
the way for relative risk aversion $\mu\in\{3,5\}$. Every obstacle it met traces to the fixed point
of the infinite-horizon problem: a condition on the response of consumption to the interest rate
fails in a tail of wealth that the stationary distribution reaches; an induction that would prove
the response bound has a step multiplier $1/\Thorn>1$, so it cannot close as a supremum argument;
and the only floor on capital supply comes from an ergodic identity that loses a factor
$\|u'\|^2/\operatorname{Var}$. Overlapping generations remove all three at once. A household
that lives $J$ periods solves a finite backward induction from an explicit terminal condition, so
a multiplier $1/\Thorn$ per step compounds to $\Thorn^{-J}$, which at Aiyagari's numbers is about
$1.03$. Wealth is bounded by $J$ periods of accumulation, so any condition needs checking on a
known bounded region. And the distribution of the youngest cohort is common to every interest
rate, so an ordering of policies pushes forward to an ordering of the whole age profile of
distributions by induction on age, with no Doeblin atom and no uniqueness of a stationary
distribution.

The target of the programme is uniqueness of the stationary equilibrium of a stochastic
life-cycle economy in the manner of \citet{Huggett1996} at $\mu\in\{3,5\}$. Section~\ref{sec:lc}
records the first stage, the life-cycle household, which is formalised. Section~\ref{sec:prog}
records the rest of the programme, the classical results of \citet{Diamond1965} included, and
what each stage needs. As in the companion, every result stated with a formal name has been
checked by the Lean kernel against Mathlib; the development is at
\url{https://github.com/LeanEconomics/LeanEconomics}, directory \texttt{LeanEconomics/OLG}.

\section{The life-cycle household}\label{sec:lc}

\subsection{Stages}

The household of the companion write-up is kept: earnings $y(z)$ follow a finite Markov chain
$\pi$, the gross return is $R=1+r$, cash on hand is $m=Ra+y(z)$, saving $a'\in[0,\bar a]$ is
bounded below by the borrowing limit and above by an artefactual cap, utility is CRRA with
$u'(c)=c^{-\mu}$ and $\beta\in(0,1)$. What changes is the horizon: the household lives $J$
periods and leaves nothing behind. Count periods from the end. With $k$ periods still to come
after today, the continuation value is the $k$-th Bellman iterate from zero,
\[
  V_k=T^k 0,\qquad (Tv)(a,z)=\max_{a'}\ \bigl\{u(m-a')+\beta\textstyle\sum_{z'}\pi(z,z')v(a',z')\bigr\},
\]
and the household with $k$ periods to come saves $g_k=\operatorname{policyOf}V_k$ and consumes
$c_k=m-g_k$ (\lean{stageValue}, \lean{stagePolicy}, \lean{stageConsumption}). Age $j$ of a
$J$-period life is stage $k=J-1-j$. In the last period the household saves nothing and consumes
its cash on hand (\lean{stagePolicy\_zero}, \lean{stageConsumption\_zero}).

This is the finite-horizon problem with a flat age profile of earnings: the same income process at
every age. An age profile, retirement and survival risk need income to be a parameter of the
Bellman step rather than a field of the structure, and are deferred to the next stage of the
programme.

\subsection{What transfers without a fixed point}

The companion development proved its household results for the Bellman operator applied to an
arbitrary continuation $v$, and only then passed to the fixed point. Here there is no fixed point
to pass to, and each result is read at $v=V_k$:

\begin{itemize}[leftmargin=1.5em]
\item consumption is positive at every stage when marginal utility is unbounded at zero
  (\lean{stageConsumption\_pos});
\item the stage value has concave slices, so the optimal action is unique and both saving and
  consumption rise with wealth (\lean{stagePolicy\_mono}, \lean{stageConsumption\_mono});
\item Carroll--Kimball: consumption is concave in wealth at every stage
  (\lean{concaveOn\_stageConsumption});
\item the Euler inequality under the cap, $\beta R\,\E[u'(c_k(a',z'))\mid z]\le u'(c_{k+1}(a,z))$,
  at every state (\lean{stage\_euler\_le}), and the Euler inequality above the floor, the reverse,
  at every state where the household saves (\lean{stage\_euler\_ge}).
\end{itemize}

\subsection{The finite-horizon sandwich}

Let $\Thorn=(\beta R)^{1/\mu}$ and define the finite-horizon minimal propensities to consume by
\[
  \kappa_0=1,\qquad \kappa_{k+1}=\frac{\kappa_kR}{\Thorn+\kappa_kR},\qquad\text{so that}\qquad
  \frac1{\kappa_k}=\sum_{i=0}^{k}\Bigl(\frac{\Thorn}{R}\Bigr)^{i}
\]
(\lean{stageMPC}): the share of cash on hand consumed by a household with $k$ periods to come and
no future income. Let $H_0=0$ and $RH_{k+1}(z)=\sum_{z'}\pi(z,z')\bigl(y(z')+H_k(z')\bigr)$, the
expected present value of the next $k$ incomes (\lean{stageHumanWealth}).

\begin{theorem}[The finite-horizon sandwich]\label{thm:sandwich}
Assume CRRA utility with marginal utility unbounded at zero, and that the cap is slack at every
stage. Then for every $k$, every $z$ and every $a\in[0,\bar a]$,
\[
  \kappa_k\,m\ \le\ c_k(a,z)\ \le\ \kappa_k\bigl(m+H_k(z)\bigr).
\]
\formal{IncomeFluctuation.stageMPC\_mul\_le\_stageConsumption,
IncomeFluctuation.stageConsumption\_le\_stageHumanWealth}
\end{theorem}

Each bound is one Euler step, applied $k$ times. The lower bound: if consumption against $v$ is at
least $\kappa m$, the Euler inequality above the floor at a state where the household saves gives
$c^{-\mu}\le\beta R\,\E[(\kappa Ra')^{-\mu}]$ with $a'=m-c$, hence $\kappa R(m-c)\le\Thorn c$ and
$c\ge\kappa R(\Thorn+\kappa R)^{-1}m$; at a corner $c=m$ (\lean{crra\_stageMPC\_step}). The upper
bound: if consumption against $v$ is at most $\kappa(m+H(z))$, the Euler inequality under the cap
and Jensen for the negative power give $c^{-\mu}\ge\beta R\,(\E[c'])^{-\mu}$ with
$\E[c']\le\kappa R(a'+H'(z))$, hence $c\le\kappa R(\Thorn+\kappa R)^{-1}(m+H'(z))$
(\lean{crra\_stageUpper\_step}). The recursion for $\kappa$ is the same on both sides, which is
why the sandwich is exact in the homogeneous case.

The propensities decrease in the horizon towards the infinite-horizon $\kappa=1-\Thorn/R$ of the
companion write-up: $\kappa\le\kappa_k$ for every $k$ (\lean{minMPC\_le\_stageMPC}), so a
household with less life left consumes a larger share of its cash, and the infinite-horizon bounds
are the limiting case. A first consequence for the programme is an explicit bound on saving,
$g_k(a,z)\le(1-\kappa_k)(Ra+y_{\max})$ (\lean{stagePolicy\_le}), which iterated from zero wealth
at the start of life bounds the wealth of every cohort.

\section{Retirement, a pension, and the means test}\label{sec:ret}

\subsection{Retirement with a lump-sum pension}

The household works until age $J_r$ and is retired after it, receiving a pension in place of its
labour earnings. With a \emph{lump-sum} pension $p$ the retired household faces the same problem
as the working one with income replaced by the constant $p$, so the retired phase is itself an
income-fluctuation problem (\lean{withPension}) and everything in Section~\ref{sec:lc} applies to
it unchanged. The life-cycle value is assembled by iterating the working Bellman operator on the
retired value: $\mathrm{lifeValue}(n_r,n_w)$ is the value of a household with $n_w$ working and
$n_r$ retired periods still to come (\lean{lifeValue}), so a household of age $j<J_r$ in a life of
$J$ periods sits at $n_w=J_r-j$ and $n_r=J-J_r$. Saving and consumption rise with wealth at every
age (\lean{lifePolicy\_mono}, \lean{lifeConsumption\_mono}).

Two facts about the pension itself are worth recording. A retiree's human wealth is the same in
every earnings state, because the pension does not depend on the earnings shock
(\lean{stageHumanWealth\_withPension\_const}), and it never exceeds the perpetuity value of the
pension, $H_k\le p/r$ (\lean{stageHumanWealth\_withPension\_le}). With Theorem~\ref{thm:sandwich}
the retiree's consumption therefore satisfies
\[
  \kappa_k\,m\ \le\ c_k(a)\ \le\ \kappa_k\Bigl(m+\frac{p}{r}\Bigr)
\]
at every retired age. And cash on hand is strictly increasing in assets
(\lean{withPension\_resources\_strictMono}). Both of these fail once the pension is means-tested.

\subsection{Means testing}

A pension is means-tested when the amount paid falls with the recipient's assets. We write the
schedule as a function $p(\cdot)$ of assets and cash on hand as $m(a)=Ra+p(a)$ (\lean{cash}),
and consider three schedules (\lean{lumpSum}, \lean{tapered}, \lean{cliff}): a lump sum; a taper,
withdrawn at rate $\tau$ per unit of assets above a threshold and never negative; and a cliff,
withdrawn entirely above the threshold.

It is often argued that a means test makes the household's policy function non-monotone, and the
reason usually given is that the test makes the value function non-concave. The argument as
usually given is wrong, and the conclusion is nevertheless true, for a different reason. Three
results separate the two.

\begin{theorem}[Monotonicity in cash on hand is never lost]\label{thm:topkis}
Let period utility be strictly concave on the positive half-line and let $W$ be \emph{any}
continuation value. If $m_1<m_2$ and $x_i$ is optimal at $m_i$, then $x_1\le x_2$.
\formal{OLG.isOptimalSaving\_mono}
\end{theorem}

The proof is Topkis. Adding the two optimality inequalities gives
$u(m_1-x_1)+u(m_2-x_2)\ge u(m_1-x_2)+u(m_2-x_1)$; if $x_2<x_1$ the two right-hand consumptions are
strict convex combinations of the two left-hand ones with the same sum, so strict concavity makes
the reverse inequality strict, a contradiction. The continuation enters only through terms that
cancel. So a non-concave value function --- which a means test certainly produces --- cannot by
itself make saving fall with cash on hand, whatever the shape of $W$.

What a non-concave continuation does produce is a jump, and the jump falls on consumption.

\begin{theorem}[Consumption falls when wealth rises]\label{thm:cons}
There are a continuation value arising from a means-tested pension and two levels of cash on hand,
the second larger, at which optimal consumption is strictly smaller. With log utility, a unit
gross return, and a pension of one withdrawn entirely above assets of one: at $m=6$ the household
saves exactly $1$ and consumes $5$; at $m=8$ it saves $4$ and consumes $4$.
\formal{OLG.exists\_consumption\_decrease}
\end{theorem}

At $m=6$ the household keeps the pension, saving the most it can without losing it; at $m=8$ the
pension is no longer worth keeping, so the household gives it up and saves the proceeds. Saving
rises from $1$ to $4$, as Theorem~\ref{thm:topkis} says it must, and since it rises by more than
wealth does, consumption falls. Both optima are exact: the objective is $\log$ of a quadratic on
each side of the threshold, and the comparison is checked by the kernel.

The third result is the genuine non-monotonicity of the saving policy, and it comes from the
budget rather than from the value function. A retiree's pension is assessed each period on the
assets then held, so the schedule enters cash on hand directly.

\begin{theorem}[A gentle taper is harmless; a cliff is not]\label{thm:taper}
If $0\le\tau<R$ then $a\mapsto Ra+p(a)$ is strictly increasing under the taper
(\lean{cash\_strictMono\_tapered}), and so, by Theorem~\ref{thm:topkis}, saving is monotone in
assets (\lean{isOptimalSaving\_mono\_assets}). Under a cliff it is not: if $a\le\bar a<b$ and
$R(b-a)<p$ then cash on hand is strictly smaller at $b$ (\lean{cash\_lt\_cash\_of\_cliff}), and
there are two asset levels, the second larger, at which the household saves strictly less: with a
concave continuation, at assets $1$ it saves $1/2$ and at assets $3/2$ only $1/4$.
\formal{OLG.exists\_saving\_decrease\_in\_assets}
\end{theorem}

The continuation in the second half of Theorem~\ref{thm:taper} is $\log(x+1)$, which is concave:
nothing is hidden in it. The non-monotonicity is entirely the budget's, and the threshold is
sharp, $\tau$ against $R$. A taper gentler than the gross return leaves the household's saving
monotone in assets no matter how badly the test distorts the value function; a taper steeper than
the gross return, or a cliff, makes a richer household poorer in the only sense that matters for
today's decision.

\section{The programme}\label{sec:prog}

The stages, in the order we intend to take them. The first is done (Section~\ref{sec:lc}).

\paragraph{Equilibrium.} A stationary overlapping-generations equilibrium: the distribution
$\mu_1$ of the youngest cohort is given (zero wealth, earnings drawn from the invariant law),
$\mu_{j+1}$ is the pushforward of $\mu_j$ under the age-$j$ policy and the earnings chain,
capital supply is $K(r)=\sum_j\omega_j\int a\,d\mu_j$ with cohort weights $\omega_j$, labour is
exogenous, and the firm is Aiyagari's. Targets: continuity of $K$ in the rate, a finite composition
of continuous pushforwards; existence by the intermediate value theorem against the firm's demand,
with the floor at the top of the interval from explicit accumulation over the life rather than from
an ergodic identity; Light's Theorem 2 by induction on age, an ordering of policies at every age
giving first-order stochastic dominance of every $\mu_j$ and so of $K$; and single crossing
against a decreasing demand.

\paragraph{Theorem 1 by backward induction.} The research core: saving rises with the rate at
every age, for $\mu\in\{3,5\}$. The reduction of the companion write-up applies stage by stage,
Theorem 1 at stage $k+1$ following from the slack elasticity condition on the stage-$k$
consumption functions at the pivot; the transfer lemma propagates an affine bound on the response
of consumption to the rate backward from the terminal condition $c_0=m$, where the response is
exactly $\Delta r\cdot a$; and the slack condition at each stage is a finite check on the
constants so obtained. The obstruction of the infinite horizon, a step multiplier of at least
$1/\Thorn$, does not bite in $J$ steps. The first session is numerical: a life-cycle economy with
$J=60$, Aiyagari's earnings chain and $\mu\in\{3,5\}$, checking the slack condition at every age
and wealth and running the recursion for the constants. If it holds the argument is formalised;
if it fails only at moderate ages and wealth we look for the age-dependent refinement; if it fails
badly, log utility is the fallback, where Theorem 1 at each stage follows from the rescaling proof
of the companion, which is already an induction on Bellman iterates.

\paragraph{Diamond (1965) and Samuelson.} The two-period economy without risk: a unique steady
state and monotone global convergence for logarithmic utility and Cobb--Douglas production,
existence under Inada conditions and a multiplicity witness in general, dynamic inefficiency of a
steady state with $r<n$ by an explicit Pareto-improving transfer, the golden rule, and the
autarkic and monetary steady states of the pure-exchange economy. These are standalone results
and will be the classical sections of this write-up.

\paragraph{Extensions.} An age profile of earnings, retirement with a pension and a government
budget, survival risk with accidental bequests as a second fixed point, population growth, and
Lasry--Lions with age as time, whose finite-horizon core is already formalised and whose
monotonicity is again the elasticity condition.


\bibliographystyle{ecta}
\bibliography{refs}

\appendix
\section{Index of formal results}\label{app:index}

All in the namespace \lean{LeanEconomics.IncomeFluctuation}, directory \lean{LeanEconomics/OLG}.
Every result depends only on the three standard axioms (\lean{propext}, \lean{Classical.choice},
\lean{Quot.sound}) and none uses \lean{sorry}.

\begingroup\footnotesize
\begin{longtable}{@{}p{0.38\textwidth}p{0.40\textwidth}p{0.16\textwidth}@{}}
\toprule
Result & Lean name & Module \\
\midrule
\endhead
Stage value, policy, consumption & \lean{stageValue}, \lean{stagePolicy}, \lean{stageConsumption} & OLG/\allowbreak LifeCycle\\
Last period: no saving & \lean{stagePolicy\_zero} & OLG/\allowbreak LifeCycle\\
Last period: consume cash on hand & \lean{stageConsumption\_zero} & OLG/\allowbreak LifeCycle\\
Positive consumption at every stage & \lean{stageConsumption\_pos} & OLG/\allowbreak LifeCycle\\
Saving and consumption rise with wealth & \lean{stagePolicy\_mono}, \lean{stageConsumption\_mono} & OLG/\allowbreak LifeCycle\\
Carroll--Kimball at every stage & \lean{concaveOn\_stageConsumption} & OLG/\allowbreak LifeCycle\\
Euler under the cap, by stage & \lean{stage\_euler\_le} & OLG/\allowbreak LifeCycle\\
Euler above the floor, by stage & \lean{stage\_euler\_ge} & OLG/\allowbreak LifeCycle\\
Finite-horizon propensities, human wealth & \lean{stageMPC}, \lean{stageHumanWealth} & OLG/\allowbreak LifeCycle\\
Propensities in $(0,1]$, above $\kappa$ & \lean{stageMPC\_pos}, \lean{stageMPC\_le\_one}, \lean{minMPC\_le\_stageMPC} & OLG/\allowbreak LifeCycle\\
Lower Euler step & \lean{crra\_stageMPC\_step} & OLG/\allowbreak LifeCycle\\
Upper Euler step & \lean{crra\_stageUpper\_step} & OLG/\allowbreak LifeCycle\\
Theorem~\ref{thm:sandwich}, lower & \lean{stageMPC\_mul\_le\_stageConsumption} & OLG/\allowbreak LifeCycle\\
Theorem~\ref{thm:sandwich}, upper & \lean{stageConsumption\_le\_stageHumanWealth} & OLG/\allowbreak LifeCycle\\
One-step saving bound & \lean{stagePolicy\_le} & OLG/\allowbreak LifeCycle\\
\bottomrule
\end{longtable}
\endgroup


\end{document}
